\scanpage{25}
\mspage{11}
est unique, et $g_i$ défini \unsure{à} multiplication à gauche par un él. quelconque de $L_i$, soit $l_i^{\alpha_i}$ $(\alpha_i\in\hat{\mathbb Z})$. On passe de $u$ à $\operatorname{int}(g)\circ u$ en \inserted{gardant $p$, et en} remplaçant les $g_i$ par les $gg_i$.

Si $v$ est donné par $q,(h_i)$, on a
\[
vu(l_i)=v(\operatorname{int}(g_i^{\deleted{\unsure{-1}}})l_i^p)
=\operatorname{int}(v(g_i^{\deleted{\unsure{-1}}}))v(l_i)^p
\]
où
\[
v(l_i)=\operatorname{int}(h_i^{\deleted{\unsure{-1}}})l_i^q
\qquad\text{donc}\qquad v(l_i)^p=\operatorname{int}(h_i^{\deleted{\unsure{-1}}})l_i^{pq}
\]
donc
\[
(23)\quad vu(l_i)=\operatorname{int}\bigl(
\overbrace{v(g_i^{\deleted{\unsure{-1}}})h_i^{\deleted{\unsure{-1}}}}^{\deleted{\unsure{(h_i v(g_i))^{-1}}}}
\bigr)l_i^{pq}.
\]
Si on désigne par $S\widehat{\widehat{\mathcal E}}^!$ l'extension de $\widehat{\widehat{\mathcal E}}^!$ par $\hat{\mathbb Z}^{\{0,1,\infty\}}$ \deleted{\unsure{définie par le groupe}} \unsure{formé} des \deleted{couples} triples d'un $u\in\widehat{\widehat{\mathcal E}}^!$ et d'un système de $g_i$ qui explicitent \inserted{via (21)} l'isomorphisme de la structure à lacets ($p$ étant le multiplicateur), on a sur $S\widehat{\widehat{\mathcal E}}^!$ trois fonctions à valeurs dans $\deleted{\unsure{\mathcal E}}\hat\pi$
\[
(u,(g_i))\overset{\gamma_i}{\longmapsto}g_i,
\]
et on aura
\[
(24)\quad\gamma_i(\tilde v\tilde u)
=\deleted{\unsure{\tilde v(\tilde u_i(\tilde u))\cdot\gamma_i(\tilde v)}}
\inserted{\deleted{\unsure{\gamma_i(\tilde v)\cdot v(\gamma_i(\tilde u))}}
\unsure{\tilde v(\gamma_i(\tilde u)\gamma_i(\tilde v)|\tilde v(\gamma_i(\tilde u))\gamma_i(\tilde v)}}
\]
i.e. $(\gamma_i)$ est un 1-cocycle sur le groupe \deleted{$\unsure{\widehat{\widehat{\mathcal E}}^!}$}\inserted{$S\widehat{\widehat{\mathcal E}}^!$} à valeurs dans le groupe \deleted{\unsure{produit du groupe}}\inserted{opposé au groupe} $\pi^{\{0,1,\infty\}}$.
\margin{\uline{NB} $S\widehat{\widehat{\mathcal E}}^!$ a été construit comme le groupe des Systèmes
\[
(p,g_0,g_1,g_\infty)\in\hat{\mathbb Z}^*\times\pi^{\{0,1,\infty\}}
\]
satisfaisant la relation (22).}
\deleted{groupe}
Le sous-groupe de $S\widehat{\widehat{\mathcal E}}^!$ noyau de $S\widehat{\widehat{\mathcal E}}^!\longrightarrow\operatorname{Aut}\operatorname{ext}_{\mathrm{lac}}^!(\hat\pi)$ s'identifie au groupe produit
\[
\left(\hat\pi\times\prod_{i\in\{0,1,\infty\}}L_i\right)
\simeq\hat\pi\times\unsure{(\hat{\mathbb Z})^{\{0,1,\infty\}}|(\hat{\mathbb Z})^3}
\]
\mspage{24}
